% ==============================================================================
% ПАРТИЯ B03-042: Список литературы: Записи 71..140 (блоки DISS-005-B2152 .. DISS-005-B2221)
% Точный перенос рукописи v0.5 без изменений
% ==============================================================================

\begin{enumerate}[resume, label=\arabic*.]
% DISS-005-B2152
\item [18] H. Zhang and S. Li, “Design of adaptive line protection under smart grid,” in 2011 International Conference on Advanced Power SystemAutomation and Protection, vol. 1. IEEE, 2011, pp. 599–603.
% DISS-005-B2153
\item [19] S. Wu, “An adaptive limited wide area differential protection for powergrid with micro-sources,” Protection and Control of Modern Power Systems, vol. 2, no. 3, pp. 1–9, 2017.
% DISS-005-B2154
\item [19] S. Wu, “An adaptive limited wide area differential protection for powergrid with micro-sources,” Protection and Control of Modern Power Systems, vol. 2, no. 3, pp. 1–9, 2017.
% DISS-005-B2155
\item [21] A. Summers, T. Patel, R. Matthews, and M. J. Reno, “Prediction of relay settings in an adaptive protection system,” in 2022 IEEE Power \& Energy Society Innovative Smart Grid Technologies Conference (ISGT). IEEE, 2022, pp. 1–5.
% DISS-005-B2156
\item [22] A. L. Kulikov, A. Loskutov, and D. Bezdushniy, “Relay protection andautomation algorithms based on machine learning,” Energies, 2022.
% DISS-005-B2157
\item [23] V. Lachugin, D. Vertoguzov, S. Danilov, V. Sazanov, A. Voloshin, andA. Kovalenko, “Development of solutions for automatic computation of relay protection and automation operation parameters,” in 2023 6th International Scientific and Technical Conference on Relay Protection and Automation (RPA), vol. 1. IEEE, 2023, pp. 1–14.
% DISS-005-B2158
\item [24] Z. e. a. Liu, “Construction of the relay protection device model data center,” IEEE Conference Publication, 2024.
% DISS-005-B2159
\item [25] A. N. Sheta, G. M. Abdulsalam, B. E. Sedhom, and A. A. Eladl, “Comparative framework for ac-microgrid protection schemes: challenges,solutions, real applications, and future trends,” Protection and control of modern power systems, vol. 8, no. 2, pp. 1–40, 2023.
% DISS-005-B2160
\item [26] H. e. a. Wang, “A setting optimization ensemble for a distributed power grid protective relay system,” Applied Sciences, vol. 14, no. 6, p. 2278, 2024.
% DISS-005-B2161
\item [27] P. Singh, U. Kumar, N. K. Choudhary, and N. Singh, “Advancements in protection coordination of microgrids: a comprehensive review of protection  challenges and mitigation schemes for grid stability,” Protection and Control of Modern Power Systems, vol. 9, no. 6, pp. 156–183, 2024.
% DISS-005-B2162
\item [28] Z. Yang, H. Wang, W. Liao, C. L. Bak, and Z. Chen, “Protection challenges and solutions for ac systems with renewable energy sources: A review,” Protection and Control of Modern Power Systems, vol. 10, no. 1, pp. 18–39, 2024.
% DISS-005-B2163
\item [29] J. W. Yan Li et al., “Fast searching of extreme operating conditions for relay protection setting calculation based on graph neural network and reinforcement learning,” arXiv preprint, 2025. [Online]. Available: \url{https://arxiv.org/abs/2501.09399}
% DISS-005-B2164
\item [30] D. A. Stepanova and A. V. Naumov, “Deep learning in relay protection of digital power industry,” International Youth Scientific and Technical Conference on Relay Protection and Automation (RPA), 2019. [Online]. Available: \url{https://www.semanticscholar.org/paper/} Deep-Learning-in-Relay-Protection-of-Digital- Power-Stepanova-Naumov
% DISS-005-B2165
\item [31] X. e. a. Li, “Fault diagnosis of intelligent substation relay protection system based on transformer architecture and migration training model,” Energy Informatics, vol. 7, p. 120, 2024.
% DISS-005-B2166
\item [32] A. Varbella et al., “Powergraph: A power grid benchmark dataset for graph neural networks,” arXiv preprint, 2024.
% DISS-005-B2167
\item [33] A. Kulikov, A. Loskutov, and D. Bezdushniy, “Relay protection and automation algorithms of electrical networks based on simulation and machine learning methods,” Energies, vol. 15, no. 18, p. 6525, 2022.
% DISS-005-B2168
\item [34] A. Kulikov, A. Loskutov, D. Bezdushniy, and I. Petrov, “Decision tree models and machine learning algorithms in the fault recognition on power lines with branches,” Energies, vol. 16, no. 14, p. 5563, 2023.
% DISS-005-B2169
\item [35] R. Machlev, A. Chachkes, J. Belikov, Y. Beck, and Y. Levron, “Open source dataset generator for power quality disturbances with deeplearning reference classifiers,” Electric Power Systems Research, vol. 195, p. 107152, 2021.
% DISS-005-B2170
\item [36] R. A. de Oliveira and M. H. Bollen, “Deep learning for power quality,”Electric Power Systems Research, vol. 214, p. 108887, 2023.
% DISS-005-B2171
\item [37] V. Pozdnyakov, A. Kovalenko, I. Makarov, M. Drobyshevskiy, and K. Lukyanov, “Adversarial attacks and defenses in fault detection and diagnosis: A comprehensive benchmark on the tennessee eastman process,” IEEE Open Journal of the Industrial Electronics Society, 2024.
% DISS-005-B2172
\item [38] H. Aguinis, N. S. Hill, and J. R. Bailey, “Best practices in data collection and preparation: Recommendations for reviewers, editors, and authors,” Organizational Research Methods, vol. 24, no. 4, pp. 678–693, 2021.
% DISS-005-B2173
\item [39] Q. Huang and T. Zhao, “Data collection and labeling techniques for machine learning,” arXiv preprint arXiv:2407.12793, 2024.
% DISS-005-B2174
\item [40] S. E. Whang, Y. Roh, H. Song, and J.-G. Lee, “Data collection and quality challenges in deep learning: A data-centric ai perspective,” The VLDB Journal, vol. 32, no. 4, pp. 791–813, 2023.
% DISS-005-B2175
\item [41] S. A. Mazhar, R. Anjum, A. I. Anwar, and A. A. Khan, “Methods of data collection: A fundamental tool of research,” Journal of Integrated Community Health (ISSN 2319-9113), vol. 10, no. 1, pp. 6–10, 2021.
% DISS-005-B2176
\item [42] A. Narayanan and E. Felten, “No silver bullet: De-identification still hard,” in IEEE Symposium on Security and Privacy, 2014, pp. 345–358.
% DISS-005-B2177
\item [43] L. Sweeney, “k-anonymity: A model for protecting privacy,” International Journal of Uncertainty, Fuzziness and Knowledge-Based Systems, vol. 10, no. 5, pp. 557–570, 2002.
% DISS-005-B2178
\item [44] “General data protection regulation (gdpr),” Official Journal of the European Union, 2016.
% DISS-005-B2179
\item [45] A. Evdakov, A. Yablokov, and G. Filatova, “Investigation of the use of rogowski coils for power direction relays in fast automatic bus transfer,” in 2024 International Conference on Industrial Engineering, Applications and Manufacturing (ICIEAM). IEEE, 2024, pp. 124–129.
% DISS-005-B2180
\item [46] G. Filatova, D. Piskunov, and A. Evdakov, “Automated method for the research of the relay protection function implemented in 6 (10) kv digital instrument transformers using real oscillograms,” in 2024 International Russian Automation Conference (RusAutoCon). IEEE, 2024, pp. 362– 367.
% DISS-005-B2181
\item [47] Y. LeCun, L. Bottou, Y. Bengio, and P. Haffner, “Gradient-based learning applied to document recognition,” Proceedings of the IEEE, vol. 86, no. 11, pp. 2278–2324, 1998.
% DISS-005-B2182
\item [48] H. Ismail Fawaz, G. Forestier, J. Weber, L. Idoumghar, and P.-A. Muller, “Deep learning for time series classification: a review,” Data mining and knowledge discovery, vol. 33, no. 4, pp. 917–963, 2019.
% DISS-005-B2183
\item [49] I. Makarov, A. Korinevskaya, and V. Aliev, “Sparse depth map interpolation using deep convolutional neural networks,” in 2018 41st international conference on telecommunications and signal processing (TSP). IEEE, 2018, pp. 1–5.
% DISS-005-B2184
\item [50] D. Chen, Y. Liu, X. Wang, and R. Zhang, “Spectral temporal graph neural network for multivariate time-series forecasting,” Advances in Neural Information Processing Systems, vol. 34, pp. 17 766–17 779, 2021.
% DISS-005-B2185
\item [51] Z. Wang, W. Yan, and T. Oates, “Time series classification using multi-channels deep convolutional neural networks,” IEEE International Conference on Data Mining, pp. 787–796, 2017.
% DISS-005-B2186
\item [52] N. V. Chawla, K. W. Bowyer, L. O. Hall, and W. P. Kegelmeyer, “Smote: Synthetic minority over-sampling technique,” Journal of Artificial Intelligence Research, vol. 16, pp. 321–357, 2002.
% DISS-005-B2187
\item [53] C. van Rijsbergen, “Information retrieval,” Butterworth-Heinemann, 1979.
% DISS-005-B2188
\item [54] A. Tharwat, “Classification assessment methods,” Applied Computing and Informatics, vol. 17, no. 1, pp. 168–192, 2020.
% DISS-005-B2189
\item [55] IEC 61850-7-4: Communication networks and systems for power utility automation, International Electrotechnical Commission Std., 2010.
% DISS-005-B2190
\item в
% DISS-005-B2191
\item В
% DISS-005-B2192
\item В
% DISS-005-B2193
\item В
% DISS-005-B2194
\item В
% DISS-005-B2195
\item В
% DISS-005-B2196
\item В
% DISS-005-B2197
\item в
% DISS-005-B2198
\item Вв
% DISS-005-B2199
\item Шабад М. А. Расчеты релейной защиты и автоматики распределительных сетей: Монография / М. А. Шабад. - 4-е изд., перераб. и доп. - СПб.: ПЭИПК, 2003 - 350 стр., ил.
% DISS-005-B2200
\item СТО 56947007-29.120.70.305-2020 Методические указания для выбора параметров настройки и срабатывания МП устройств РЗА оборудования 6-35 кВ объектов ЕНЭС (Дата введения: 15.05.2020), ПАО «ФСК ЕЭС».
% DISS-005-B2201
\item Никулов И.И., Жуков В.А., Пупин В.М.. Комплекс БАВР. Быстродействие повышает надежность электроснабжения // Новости ЭлектроТехники № 4(76) 2012. С. 2–4.
% DISS-005-B2202
\item Сульчакова А.Ю. Испытания промышленных образцов устройств быстродействующего автоматического ввода резерва (БАВР) для нефтеперекачивающих станций на RTDS. Разработка тестовой схемы для испытаний устройств БАВР на RTDS: докл. на заседании АО «НТЦ ЕЭС» 2020 г.
% DISS-005-B2203
\item Андреенко Ю. А. Измерительное устройство дифференциальной токовой защиты шин / Ю. А. Андреенко [и др.]. - (Электроэнергетика). - Текст : непосредственный // Известия вузов. Проблемы энергетики. - 2011. - № 7/8. - С. 62-70 : схемы, граф. - Библиогр.: с. 69-70 (10 назв. ). - Примеч.: с. 70. - ISSN 1998-9903.
% DISS-005-B2204
\item Siamak Bahari. A New Stabilizing Method of Differential Protection Against Current Transformer Saturation Using Current Derivatives / Siamak Bahari, Tohid Hasani, Heresh Seyedi // 14th International Conference on Protection \& Automation in Power System - pp. 33-38, April 2020
% DISS-005-B2205
\item Chothani, N.G., Bhalja, B.R., New algorithm for current transformer saturation detection and compensation based on derivatives of secondary currents and Newton’s backward difference formulae, IET Generation, Transmission \& Distribution, Vol. 8, pp. 841-850, May 2014
% DISS-005-B2206
\item Руководство по эксплуатации КРУ TER\_Sec10\_Etalon\_T1 / Версия 5.0 – 102 стр., ил.
% DISS-005-B2207
\item Руководство по эксплуатации. Руководство по эксплуатации. Общие технические требования. Комплектное устройство ШЭ-АПС-БАВР шкаф автоматики быстродействующего автоматического ввода резерва в сетях 6÷35 кв / АПДЛ.656435003 РЭ ред. 3 от 31.08.2023 – 44 стр., ил.
% DISS-005-B2208
\item Руководство по эксплуатации. Руководство оператора. Микропроцессорное устройтсво быстродействующего автоматического ввода резерва с восстановлением нормального режима электроснабжения БАВР «МИР» / АПДЛ.656121003 РЭ2 ред. 1 от 13.07.2023 – 179 стр., ил.
% DISS-005-B2209
\item Ломтев Е.А. Принципы построения и конструкции высоковольных ёмкостных делителей напряжения / Е.А. Ломтев, Д.И. Нефедьев // Измерение. Мониторинг. Управление. Контроль. – 2014 - №4(10). – с 4-8
% DISS-005-B2210
\item Roman Hrbac. A Development of a Capacitive Voltage Divider for High Voltage Measurement as Part of a Combined Current and Voltage Sensor / Roman Hrbac, Vaclav Kolar, Mikolaj Bartlomiejczyk, Tomas Mlcak// ELEKTRONIKA IR ELEKTROTECHNIKA, ISSN 1392-1215, VOL. 26, NO. 4, 2020 - pp. 25-31
% DISS-005-B2211
\item Tobias Kuhnke. Frequency Analysis of a novel 10 kV RC Voltage Divider Module composed of SMD-Components on Printed Circuit Boards / Tobias Kuhnke; Saskia Düsdieker; Frank Jenau; Mohammad Razavifar // IEEE Belgrade PowerTech, August 2023 – pp ХХ-ХХ
% DISS-005-B2212
\item You Li. Study on methods of reduce resistive voltage dividers' phase shift / You Li; Li Yalu; Liu Min; Jin Haibin // IEEE 2011 10th International Conference on Electronic Measurement \& Instruments, October 2011 – pp. 90-93
% DISS-005-B2213
\item Гречухин В.Н. Электронные трансформаторы тока и напряжения.состояние, перспективы развития и внедрения на ОРУ 110–750 кв станций и подстанций энергосистем / Гречухин В.Н. // «Вестник ИГЭУ» Вып. 4 2006 г., стр. 1-9
% DISS-005-B2214
\item Клименко К.А. Исследование электромагнитного поля трансформаторов промышленной частоты с короткозамкнутыми витками: дис. ... канд. тех. наук: 05.09.01: защищена 25.12.13: утв. 05.05.12 / Клименко Ксения Александровна., 2013. - 157 с.
% DISS-005-B2215
\item Кужеков С.Л. О проблеме выбора и замены трансформаторов тока для устройств релейной защиты / С.Л. Кужеков, А.А. Дегтярёв, Н.А. Дони, А.А. Шурупов// Известия вузов. Электромеханика. 2020. Т. 63. № 6 - с 72-82
% DISS-005-B2216
\item PC37.235/D15, Mar 2021 - IEEE Approved Draft Guide for the Application of Rogowski Coils Used for Protective Relaying Purposes, November 2021 - pp XXX
% DISS-005-B2217
\item Zhang Zhou sheng. The analysis of the factors affecting the performance of Clamp Rogowski coil / Zhang Zhou sheng, Bao Han, Ma ai qing // 2014 International Conference on Power System Technology (POWERCON 2014) – pp. 1446-1450
% DISS-005-B2218
\item Esa-Pekka Suomalainen. Onsite Calibration of a Current Transformer Using a Rogowski Coil / Esa-Pekka Suomalainen, Jari K. Hällström // IEEE TRANSACTIONS ON INSTRUMENTATION AND MEASUREMENT, VOL. 58, NO. 4, APRIL 2009 – pp. 1054-1058
% DISS-005-B2219
\item Peng Wang. Planar Rogowski Coil Current Transducer Used for Three-phase Plate-form Busbars / Peng Wang, Guixin Zhang, Xiaomei Zhu, Chengmu Luo // Instrumentation and Measurement Technology Conference - IMTC 2007 Warsaw, Poland, May 1-3, 2007 – pp. 1-4
% DISS-005-B2220
\item Lj. A. Kojovic. Comparative Performance Characteristics of Current Transformers and Rogowski Coils used for Protective Relaying Purposes / Lj. A. Kojovic // IEEE Power Engineering Society General Meeting - July 2007, pp. 1-6
\end{enumerate}

% DISS-005-B2221
\par\vspace{0.2cm}
